% ==============================================================================
% SEÇÃO 3: FASES DO PROJETO (MODELAGEM E EXPERIMENTOS)
% Autores: Gustavo (#2), Alêkson (#3), Italo (#5), Danilo (#7)
% 
% REGRA DE OURO ANTI-REDUNDÂNCIA:
% Não descrever telas ou botões do Streamlit nesta seção (a interface já está 
% mapeada na Seção 2). Focar estritamente em formulação matemática, engenharia 
% de atributos e resultados empíricos.
% ==============================================================================

\section{Fases do Projeto: Modelagem e Experimentos}
\label{sec:fases_do_projeto}

Este capítulo detalha a fundamentação teórica, a modelagem matemática e os experimentos computacionais conduzidos ao longo das três fases do ciclo de desenvolvimento do recomendador.

% ------------------------------------------------------------------------------
% 3.1 FASE 1: EDA E AS 5 HIPÓTESES DE MERCADO
% Autor: Gustavo Oliveira da Silva (Issue #2)
% ------------------------------------------------------------------------------
\subsection{Fase 1.1: Análise Exploratória de Dados (EDA) e Hipóteses de Mercado}
\label{sec:fase1_eda}

% [PLACEHOLDER - GUSTAVO: Inserir a análise das 5 hipóteses mercadológicas]
% - Contextualização do dataset Kaggle (123.849 vagas);
% - H1 a H5 com gráficos de distribuição salarial e taxas de conversão (CTR);
% - Apresentação das 3 confirmadas e 2 refutadas com prudência epistêmica;
% - Justificativa de uso do CTR como bônus de engajamento no recomendador.
\noindent\textit{[Espaço reservado para a redação da Seção 3.1 por Gustavo Oliveira -- Issue \#2. Foco nos achados estatísticos e gráficos das 5 hipóteses.]}

% ------------------------------------------------------------------------------
% 3.2 FASE 1: FILTRAGEM BASEADA EM CONTEÚDO (CBF)
% Autor: Alêkson Castro (Issue #3)
% ------------------------------------------------------------------------------
\subsection{Fase 1.2: Filtragem Baseada em Conteúdo (CBF)}
\label{sec:fase1_cbf}

% [PLACEHOLDER - ALÊKSON: Inserir a engenharia semântica da CBF]
% - Identificador textual composto: Texto = (Título x 2) + Skills + Senioridade;
% - Matriz TF-IDF com 3.000 termos e vetor do perfil do usuário;
% - Cálculo formal da similaridade do cosseno;
% - Avaliação quantitativa da CBF sem bônus (Precision@10 = 0,882 e NDCG@10 = 0,750).
\noindent\textit{[Espaço reservado para a redação da Seção 3.2 por Alêkson Castro -- Issue \#3. Foco na formulação do TF-IDF, cosseno e avaliação quantitativa offline.]}

\paragraph{Formulação Matemática e Engenharia de Atributos (Conteúdo-Semente):}
O pipeline analítico da filtragem por conteúdo apoia-se em quatro formulações centrais:

\begin{enumerate}[label=\alph*)]
    \item \textbf{Construção do Identificador Textual Composto:} O documento representativo de cada vaga é construído com peso dobrado atribuído ao título da posição:
    \begin{equation}
        \text{Texto}_j = (\text{Título}_j \times 2) \oplus \text{Skills}_j \oplus \text{Senioridade}_j
    \end{equation}
    \item \textbf{Vetorização Esparsa via TF-IDF:} Mapeamento do catálogo em matriz esparsa contendo os 3.000 termos de maior frequência inversa no documento.
    \item \textbf{Composição Dinâmica do Vetor do Perfil ($\mathbf{u}$):} Síntese aditiva a partir das seleções positivas ($\alpha$), seleções negativas ($\beta$) e inserções em texto livre ($\gamma$):
    \begin{equation}
        \mathbf{u} = \alpha \cdot \mathbf{v}_{\text{positivas}} - \beta \cdot \mathbf{v}_{\text{negativas}} + \gamma \cdot \mathbf{v}_{\text{texto}}
    \end{equation}
    \item \textbf{Similaridade do Cosseno:} Mensuração do alinhamento semântico vetorial entre o perfil sintetizado $\mathbf{u}$ e cada vaga $\mathbf{d}_j$:
    \begin{equation}
        \cos(\mathbf{u}, \mathbf{d}_j) = \frac{\mathbf{u} \cdot \mathbf{d}_j}{\|\mathbf{u}\| \|\mathbf{d}_j\|}
    \end{equation}
    \item \textbf{Desempate via Bônus Empírico de Atratividade (CTR):} Incorporação do CTR mediano obtido na Fase 1.1 para priorizar posições com apelo comprovado de mercado:
    \begin{equation}
        \text{Score Final}_j = \cos(\mathbf{u}, \mathbf{d}_j) + \alpha_{\text{CTR}} \cdot \text{CTR}_j
    \end{equation}
\end{enumerate}

% ------------------------------------------------------------------------------
% 3.3 FASE 2: FILTRAGEM COLABORATIVA (SVD VS. KNN)
% Autor: Italo Emannuel Beckman (Issue #5)
% ------------------------------------------------------------------------------
\subsection{Fase 2: Filtragem Colaborativa (Fatores Latentes SVD vs. KNN)}
\label{sec:fase2_cf}

% [PLACEHOLDER - ITALO: Inserir o confronto algorítmico e auditoria]
% - Matriz esparsa de 602k ratings sintéticos (5.000 personas x 6.000 vagas);
% - Formulação matemática aditiva do SVD (k=20) vs. KNN centrado;
% - Auditoria de métricas: RMSE SVD (0,625) vs. KNN (1,263);
% - Proximidade de 1,20x em relação ao piso do oráculo (0,5214);
% - Teste pareado de Wilcoxon com p ≈ 0 comprovando significância estatística.
\noindent\textit{[Espaço reservado para a redação da Seção 3.3 por Italo Beckman -- Issue \#5. Foco na fatoração SVD, decomposição de viés e significância estatística.]}

\paragraph{Formulação Matemática do Modelo Latente SVD (Conteúdo-Semente):}
A predição das notas estimadas para pares usuário-item $(u, i)$ não observados é computada por meio da formulação aditiva clássica de Decomposição em Valores Singulares com regularização:

\begin{equation}
    \hat{y}_{u,i} = \mu + b_u + b_i + q_i^T p_u
\end{equation}
onde:
\begin{itemize}
    \item $\mu = 3{,}230$: média global das avaliações observadas na partição de treino;
    \item $b_u \in \mathbb{R}$: viés característico do candidato $u$ (tendência sistemática de rigor ou benevolência nas notas);
    \item $b_i \in \mathbb{R}$: viés intrínseco da vaga $i$ (popularidade ou atratividade geral frente a todo o catálogo);
    \item $q_i^T p_u$: produto interno dos vetores de fatores latentes em dimensão reduzida ($k=20$), capturando a sinergia e casamento de afinidade entre candidato e anúncio.
\end{itemize}

% ------------------------------------------------------------------------------
% 3.4 FASE 3: FILTRAGEM HÍBRIDA DE ROBIN BURKE
% Autor: Danilo Belém Oliveira (Issue #7)
% ------------------------------------------------------------------------------
\subsection{Fase 3: Filtragem Híbrida (Taxonomia de Robin Burke)}
\label{sec:fase3_hibrido}

% [PLACEHOLDER - DANILO: Inserir a convergência híbrida e trade-offs]
% - Fundamentação na taxonomia clássica de Robin Burke (2002);
% - Equação da fusão ponderada linear (0,4 CBF + 0,6 CF) sobre escore bruto normalizado;
% - Política de chaveamento dinâmico (Switching) para cold start (N < 5);
% - Equilíbrio entre cobertura textual imediata e serendipidade comportamental coletiva.
\noindent\textit{[Espaço reservado para a redação da Seção 3.4 por Danilo Belém -- Issue \#7. Foco nas equações de normalização min-max contínua e regras de chaveamento.]}

\paragraph{Concepção Matemática da Hibridização (Conteúdo-Semente):}
Para neutralizar a propensão à hipersensibilidade do CBF e a vulnerabilidade ao \textit{cold start} do CF, adota-se a conjugação formal de duas técnicas da taxonomia de Robin Burke~\cite{burke2002}:

\begin{tcolorbox}[burkebox, title={Concepção Matemática da Filtragem Híbrida}]
\textbf{1. Estratégia Ponderada (\textit{Weighted}):}
\begin{equation}
    \text{Score Híbrido}_{u,j} = w_1 \cdot \text{Score}_{\text{CBF}}^{\text{norm}}(u, j) + w_2 \cdot \text{Score}_{\text{CF}}^{\text{norm}}(u, j)
\end{equation}
onde os escores semântico e colaborativo são normalizados para o intervalo comum $[0, 1]$ e submetidos à calibração paramétrica ponderada (e.g., $w_1 = 0{,}40$ para relevância de competências e $w_2 = 0{,}60$ para atratividade comportamental coletiva).

\vspace{0.25cm}
\textbf{2. Estratégia de Chaveamento para Cold Start (\textit{Switching}):}
\begin{equation}
    \text{Regime de Operação}(u) = 
    \begin{cases} 
        \text{CBF (Conteúdo Exclusivo)}, & \text{se } |I_u| = 0 \text{ (Usuário Novo)}, \\
        \text{Fusão Ponderada (CBF} + \text{CF)}, & \text{se } |I_u| \ge 1 \text{ (Usuário Ativo)},
    \end{cases}
\end{equation}
onde $|I_u|$ quantifica o total de interações históricas acumuladas pelo usuário $u$.
\end{tcolorbox}
